\subsection{Outer measure of an arbitrary set}

DEFINITION 4. --- Let \(\mu\) be a positive measure on X; for every subset A of X, the upper integral \(\mu^{*}(\varphi_{\mathrm{A}})\) is called the outer measure of A (with respect to the measure \(\mu\) ) and is denoted \(\mu^{*}(\mathrm{A})\) .

The outer measure of a set is thus a number \(\geqslant0\) , finite or equal to \(+\infty\) , that, for an open set, coincides with the outer measure defined in Def. 2 of No.~2.

PROPOSITION 16. --- If A and B are two subsets of X such that A ⊂ B, then \(\mu^{*}(A) \leqslant \mu^{*}(B)\) .

COROLLARY. --- Every relatively compact set in X has finite outer measure.

For, such a set is contained in a relatively compact open set (GT, I, §9, No.~7, Prop. 10), whose outer measure is finite (No.~2, Prop. 5).

PROPOSITION 17. --- If \((\mathrm{A}_{n})\) is an increasing sequence of subsets of X, then \(\mu^{*}(\bigcup_{n}\mathrm{A}_{n})=\sup_{n}\mu^{*}(\mathrm{A}_{n})\) .

PROPOSITION 18. --- For every sequence (A \(_{n}\) ) of subsets of X,

\[
\mu^ {*} \left(\bigcup_ {n} \mathrm{A} _ {n}\right) \leqslant \sum_ {n} \mu^ {*} (\mathrm{A} _ {n}).
\]

These propositions are the translations of Props. 10 and 13 and of Th. 3 of No.~3 for characteristic functions of sets.

PROPOSITION 19. --- For every subset A of X, \(\mu^{*}(A)\) is the infimum of the outer measures of the open sets containing A.

The proposition is obvious if \(\mu^{*}(\mathbf{A}) = +\infty\) . In the contrary case, for every \(\varepsilon\) such that \(0 < \varepsilon < 1\) , there exists a function \(f \in \mathcal{I}_{+}\) such that \(\varphi_{\mathrm{A}} \leqslant f\) and \(\mu^{*}(\mathrm{A}) \leqslant \mu^{*}(f) \leqslant \mu^{*}(\mathrm{A}) + \varepsilon\) . Let \(\mathbf{G}\) be the set of \(x \in \mathbf{X}\) such that \(f(x) > 1 - \varepsilon\) . Since \(f\) is lower semi-continuous, \(\mathbf{G}\) is open (GT, IV, §6, No.~2, Prop. 1) and contains A; on the other hand \(f \geqslant (1 - \varepsilon)\varphi_{\mathrm{G}}\) , whence

\[
\mu^ {*} (\mathrm{G}) \leqslant \frac {1}{1 - \varepsilon} \mu^ {*} (f) \leqslant \frac {1}{1 - \varepsilon} \left(\mu^ {*} (\mathrm{A}) + \varepsilon\right);
\]

since \(\varepsilon\) is arbitrary, we see that \(\mu^{*}(\mathbf{G})\) differs as little as we please from \(\mu^{*}(\mathbf{A})\) , whence the proposition.

